%! TeX program = lualatex
% SPDX-License-Identifier: CC-BY-4.0

%% START PLOS TEMPLATE
\documentclass[10pt,letterpaper]{article}
\usepackage[top=0.85in,left=2.75in,footskip=0.75in]{geometry}

% amsmath and amssymb packages, useful for mathematical formulas and symbols
\usepackage{amsmath,amssymb}

% Use adjustwidth environment to exceed column width (see example table in text)
\usepackage{changepage}

% textcomp package and marvosym package for additional characters
\usepackage{textcomp,marvosym}

% cite package, to clean up citations in the main text. Do not remove.
\usepackage{cite}

% Use nameref to cite supporting information files (see Supporting Information section for more info)
\usepackage{nameref,hyperref}

% line numbers
\usepackage[right]{lineno}

% ligatures disabled
\usepackage[nopatch=eqnum]{microtype}
\DisableLigatures[f]{encoding = *, family = * }

% color can be used to apply background shading to table cells only
\usepackage[table]{xcolor}

% array package and thick rules for tables
\usepackage{array}

% create "+" rule type for thick vertical lines
\newcolumntype{+}{!{\vrule width 2pt}}

% create \thickcline for thick horizontal lines of variable length
\newlength\savedwidth
\newcommand\thickcline[1]{%
  \noalign{\global\savedwidth\arrayrulewidth\global\arrayrulewidth 2pt}%
  \cline{#1}%
  \noalign{\vskip\arrayrulewidth}%
  \noalign{\global\arrayrulewidth\savedwidth}%
}

% \thickhline command for thick horizontal lines that span the table
\newcommand\thickhline{\noalign{\global\savedwidth\arrayrulewidth\global\arrayrulewidth 2pt}%
\hline
\noalign{\global\arrayrulewidth\savedwidth}}


% Remove comment for double spacing
%\usepackage{setspace}
%\doublespacing

% Text layout
\raggedright
\setlength{\parindent}{0.5cm}
\textwidth 5.25in
\textheight 8.75in

% Bold the 'Figure #' in the caption and separate it from the title/caption with a period
% Captions will be left justified
\usepackage[aboveskip=1pt,labelfont=bf,labelsep=period,justification=raggedright,singlelinecheck=off]{caption}
\renewcommand{\figurename}{Fig}

% Please use the included `plos2025.bst` as your BibTeX style
\bibliographystyle{plos2025}

% Remove brackets from numbering in List of References
\makeatletter
\renewcommand{\@biblabel}[1]{\quad#1.}
\makeatother

% Header and Footer with logo
\usepackage{lastpage,fancyhdr,graphicx}
\usepackage{epstopdf}
%\pagestyle{myheadings}
\pagestyle{fancy}
\fancyhf{}
%\setlength{\headheight}{27.023pt}
%\lhead{\includegraphics[width=2.0in]{PLOS-submission.eps}}
\rfoot{\thepage/\pageref{LastPage}}
\renewcommand{\headrulewidth}{0pt}
\renewcommand{\footrule}{\hrule height 2pt \vspace{2mm}}
\fancyheadoffset[L]{2.25in}
\fancyfootoffset[L]{2.25in}
\lfoot{\today}
%% END PLOS TEMPLATE

% Import custom style file containing common packages and options
\usepackage{preamble}

% Include definitions and custom commands
\input{custom-definitions}

% Import tikz options and preamble
\input{tikz/preamble.tex}

% Set data search path
\pgfplotsset{
    table/search path={../../data/processed},
}
\newcommand\dataSearchPath{../../data/processed}

\begin{document}
\vspace*{0.2in}

\begin{flushleft}
{\Large
\textbf{Supplementary Information for\\
``Evolutionary Kuramoto dynamics unravels origins of chimera states
in neural populations''}
}
\newline
\\
Thomas Zdyrski,
Scott Pauls,
Feng Fu
\end{flushleft}

\tableofcontents
\section{Two-player game order graphs}\label{sec:order_graphs}
\begin{figure*}
\begin{adjustwidth}{-2.25in}{0in}
	\centering
    \includestandalone{tikz/game-payoffs}
		\caption{}\label{fig:order_graphs}
\end{adjustwidth}
\end{figure*}

\Cref{fig:order_graphs} shows the
order graphs for the 16 symmetric, ordinal two-player games.
Each sub-panel shows the payoff order for both the column
and row player.
The letters correspond to the following payoff matrix
\begin{equation*}
  \begin{bNiceMatrix}[first-col,first-row,hvlines]
    & \textcolor{blue}{C} & \textcolor{blue}{N} \\
    \textcolor{red}{C} & A & B \\
    \textcolor{red}{N} & C & D
  \end{bNiceMatrix}
\end{equation*}
The horizontal axis represents the order of row player payoffs,
and the red lines represent row player options.
The vertical axis represents the order of column player payoffs,
and the blue lines represent column player options.
The arrows show the direction of increasing
payoff for the relevant player,
and solid black dots represent Nash equilibria.
This figure is adapted from the topological taxonomy
order graphs~\cite{bruns2015names}.

\section{Well-mixed communicative fraction with symmetry breaking}\label{sec:analytic_comm_frac}
Here, we follow the steady-state communicative fraction derivation
in the previous EK study~\cite{tripp2022evolutionary}
and extend it to include a payoff asymmetry $\alpha$.
To keep the derivation tractable,
we only consider the well-mixed case.
We will assume a low mutation rate
so that any mutation fixates prior to the next mutation.
We then consider a population with one communicative species $E = (C, \phi_i)$
and one non-communicative species $F = (N, \phi_j)$.
Our first goal is to calculate
the (communicative) $\rho_E$ and (non-communicative) $\rho_F$
fixation probabilities of a single $E$ or $F$ invader, respectively.
First, we use our $\alpha$-dependent payoff matrix,
to calculate the $\alpha$-dependent average payoff
of each strategy across the entire population,
\begin{align}
  \pi_E(k) &= \ab(\frac{k-1}{n-1}) \pab{B_0 - c}
                + \ab(\frac{n-k}{n-1}) \ab(\beta(\Delta\phi) 2\alpha - c) \\
           &= \frac{1}{n-1}
               \ab(k \ab(B_0 - 2 \alpha \beta(\Delta\phi))
                 + 2 \alpha n \beta(\Delta\phi) - B_0 - (n-1) c)
\end{align}
and
\begin{align}
  \pi_F(k) &= \ab(\frac{k}{n-1}) \ab(\beta(\Delta\phi) 2 (1-\alpha))
                + \ab(\frac{n-k-1}{n-1}) \ab(0) \\
           &= \ab(\frac{k}{n-1}) 2 (1-\alpha) \beta(\Delta\phi)
\end{align}
Here, $k \in [0,n]$ is the number of $E$-players
$\Delta \phi$ is the phase difference between the $E$ and $F$
strategies,
and we defined $\beta(\Delta \phi) = \beta_0 f(\Delta \phi)$.
Note that the first term of $\pi_E$ corresponds
to the $E$-$E$ payoff, which has $\Delta \phi = 0$,
hence $B_0 f(\Delta \phi) = B_0$.
As mentioned in the main text,
each node has an exponential total fitness
$f_E = \exp(\delta (n-1) \pi_E)$
or
$f_F = \exp(\delta (n-1) \pi_F)$,
depending on its strategy,
with selection strength $\delta$.

In this paragraph, we summarize the relevant steps;
see the prior work~\cite{tripp2022evolutionary}
for more detailed derivations.
We now model the Moran process
as a Markov chain where each state $i \in \Bab{0,1,\ldots,n}$
is the number of (communicative) $E$ strategies.
Interestingly, we note that since the Moran process changes, at most,
one strategy per time step,
this Markov chain has the same fixation probabilities
as a (state-dependent) ``gambler's ruin'' problem with ties.
We denote $x_i$ as the probability that state $i$ will fixate
to state $n$ with all players using strategy $E$.
Then, the $E$ and $F$ fixation probabilities defined above are given by
$\rho_E = x_1$ and $\rho_F = 1 - x_{n-1}$.
Now, the Markov chain is described by the transition probability $p_{i,j}$
from state $i$ to state $j$ as
\begin{align*}
  p_{0,0} &= 1, \\
  p_{n,n} &= 1, \\
  p_{i,i-1} &= \frac{i}{n} \frac{(n-i) f_F(i)}{i f_E(i) + (n-i) f_F(i)} \\
  p_{i,i+1} &= \frac{n-i}{n} \frac{i f_E(i)}{i f_E(i) + (n-i) f_F(i)} \\
  p_{i,i} &= 1 - p_{i,i+1} - p_{i,i-1}
\end{align*}
Then, we obtain a recurrence relation for $x_i$ by conditioning
on the outcome of the first step:
\begin{align*}
  x_i &= x_{i-1} p_{i,i-1} + x_i p_{i,i} + x_{i+1} p_{i,i+1} \\
\end{align*}
with boundary values
\begin{align*}
  x_0 &= 0 \\
  x_1 &= 1
\end{align*}
Using the $p_{i,i} = 1 - p_{i,i+1} - p_{i,i-1}$ relation from above
to replace the $p_{i,i}$ term gives
\begin{equation*}
  \pab{x_i - x_{i-1}} \gamma_i = \pab{x_{i+1} - x_i}
\end{equation*}
Then, defining $y_i \coloneqq x_i - x_{i-1}$,
we find $y_1 = x_1$ and $y_{i+1} = \gamma_i y_i$,
yielding $y_i = \prod_{j=1}^{i-1} \gamma_j x_1$ for $i \ge 2$.
Finally, we form a telescoping sum to yield
\begin{equation*}
  1 - x_1 = x_n - x_1
  = \sum_{i=1}^{n-1} y_{i+1}
  = \sum_{i=1}^{n-1} \prod_{j=1}^i \gamma_j x_1
\end{equation*}
Solving this for $x_1$ gives
\begin{equation}
  \rho_E = x_1 = \frac{1}{1 + \sum_{i=1}^{n-1} \prod_{j=1}^i \gamma_j}
  \label{eq:comm_fixation_prob}
\end{equation}
Likewise, $x_i = \sum_{j=1}^i y_i = \sum_{j=0}^{i-1} y_{i+1}$, so
\begin{equation*}
  x_i
  = x_1 + \sum_{j=1}^{i-1} \prod_{k=1}^j \gamma_k x_1
  = \frac{1 + \sum_{j=1}^{i-1} \prod_{k=1}^j \gamma_k}
    {1 + \sum_{j=1}^{n-1} \prod_{k=1}^j \gamma_k}
\end{equation*}
Therefore, we find
\begin{equation*}
  \rho_F = 1 - x_{n-1}
  = 1 - \frac{1 + \sum_{j=1}^{n-2} \prod_{k=1}^j \gamma_k}
    {1 + \sum_{j=1}^{n-1} \prod_{k=1}^j \gamma_k}
  = \frac{\prod_{k=1}^{n-1} \gamma_k}{1 + \sum_{j=1}^{n-1} \prod_{k=1}^j \gamma_k}
\end{equation*}
This implies
\begin{equation}
  \frac{\rho_F}{\rho_E} = \prod_{k=1}^{n-1} \gamma_k
  \label{eq:fixation_prob_ratio}
\end{equation}
Therefore, we've found $\rho_E$ and $\rho_F$ for a population
with two fixed strategies.

Now, we need to derive the form of the $\gamma_k$,
which will include the new asymmetry factor $\alpha$:
\begin{align*}
  \gamma_k &= \frac{f_F(k)}{f_E(k)} \\
           &= \exp[\delta (n-1) (\pi_F(k) - \pi_E(k))]
           \\
           &= \exp\ab[\delta \ab(
    k 2 (1-\alpha) \beta(\Delta\phi)
    - k \ab(B_0 - 2 \alpha \beta(\Delta\phi))
                 - 2 \alpha n \beta(\Delta\phi) + B_0 + (n-1) c)] \\
      &= \exp\ab[\delta \ab(
    \ab(2 \beta(\Delta\phi) - B_0) k
                 + B_0 - 2 \alpha n \beta(\Delta\phi) + (n-1) c)]
\end{align*}
We also calculate
\begin{equation}
  \begin{aligned}
    \prod_{k=1}^j  \gamma_k
      &= \prod_{k=1}^j \exp\ab[\delta \ab(
      \ab(2 \beta(\Delta\phi_{qr}) - B_0) k
      + B_0 - 2 \alpha n \beta(\Delta\phi_{qr}) + (n-1) c)]
      \\
      &= \exp\ab[\delta \sum_{k=1}^j \ab(
      \ab(2 \beta(\Delta\phi_{qr}) - B_0) k
      + B_0 - 2 \alpha n \beta(\Delta\phi_{qr}) + (n-1) c)]
      \\
      &= \exp\ab[\delta \ab(
      \ab(2 \beta(\Delta\phi_{qr}) - B_0) \frac{j(j+1)}{2}
      + j \ab(B_0 - 2 \alpha n \beta(\Delta\phi_{qr}) + (n-1) c))]
      \\
      &= \exp\ab[\delta \ab(
      \ab(\beta(\Delta\phi_{qr}) - \frac{B_0}{2}) j^2
      + j \ab(\frac{B_0}{2} + \beta(\Delta\phi_{qr}) (1 - 2 \alpha n)  + (n-1) c))]
  \end{aligned}
  \label{eq:gamma_prod}
\end{equation}
Additionally, we note that substituting
\cref{eq:gamma_prod} into \cref{eq:fixation_prob_ratio}
with $j=n-1$ shows that $\rho_F/\rho_E$
depends of $\Delta \phi_{qr}$.
This is in contrast to the $\alpha=0.5$ case~\cite{tripp2022evolutionary}
where $\rho_F/\rho_E$ simplifies to the $\Delta \phi_{qr}$-independent form
$\rho_F/\rho_E = \exp \ab(\delta \ab[(n-1)c - B_0 (n-2)/2])$.

Next, we consider the long-time dynamics incorporating more than two strategies.
By again using the low-mutation-rate assumption,
we can assume that any mutation either fixates
or dies out before the next mutation.
Therefore, the system evolves from one homogeneous state to another
using the probabilities we just derived.
We can model this behavior with a new Markov chain
using the $m=20$ phases and two communicative strategies (C and N)
to give a $2m$-dimensional state space:
$\Bab{(C,\phi_1), \ldots (C,\phi_m), (N,\phi_1), \ldots (N,\phi_m)}$.
Then, the transition probability $\rho_{NC,\Delta \phi_{qr}}$
from state $(N,\phi_q)$ to state $(C,\phi_r)$ is just $\rho_E$ above.
Similarly, the probability $\rho_{CN,\Delta \phi_{rq}}$
that a $(C,\phi_r)$ state is successfully invaded by
an $(N,\phi_q)$ state is simply $\rho_F$ above.
By the rotation symmetry of $\phi$, each $\phi_i$ has the same probability of fixation.
Therefore, consider an arbitrary $\phi_q$ and denote its communicative fixation probability as
$s_1$ and its non-communicative fixation probability $s_2$.
We can calculate the ratio of these probabilities
by incorporating the fixation probabilities against all other $\phi_r$ phases:
\begin{align}
  \frac{s_2}{s_1} &= \frac
    {\sum_{r=1}^{m} \rho_{CN,\Delta \phi_{rq}}}
    {\sum_{r=1}^{m} \rho_{NC,\Delta \phi_{qr}} }
  \\
  &=
  \frac
  {\sum_{r=1}^{m} \rho_{NC,\Delta \phi_{qr}}
    \exp \Bab{
      \delta (n-1)
      \bab{
        (n-1) c + n \beta(\Delta \phi_{qr}) (1 - 2 \alpha)
        - \frac{n-2}{2} B_0
      }
    }
  }
  {\sum_{r=1}^{m} \rho_{NC,\Delta \phi_{qr}}}
  \\
  &=
  \bab{
    \frac
    {\sum_{r=1}^{m} \frac
      {\exp \Bab{ \delta (n-1)
        \bab{
          (n-1) c + n \beta(\Delta \phi_{qr}) (1 - 2 \alpha) - \frac{n-2}{2} B_0
        }}
      }
      {1 + \sum_{j=1}^{n-1} \exp\Bab{
        \delta \bab{
         \pab{\beta(\Delta \phi_{qr}) - B_0/2} j^2
         + j \pab{B_0/2 + \beta(\Delta \phi_{qr}) (1 - 2 \alpha n) + (n-1) c}
        }}
      }
    }
    {\sum_{r=1}^{m} \frac{1}
      {1 + \sum_{j=1}^{n-1} \exp\Bab{
        \delta \bab{
         \pab{\beta(\Delta \phi_{qr}) - B_0/2} j^2
         + j \pab{B_0/2 + \beta(\Delta \phi_{qr}) (1 - 2 \alpha n) + (n-1) c}
        }}
      }
    }
  }
  \label{eq:full_analytic_frac}
\end{align}
The first equality used \cref{eq:fixation_prob_ratio}
and the second equality used \cref{eq:comm_fixation_prob}.
Unlike in the $\alpha = 0.5$ symmetric case~\cite{tripp2022evolutionary},
the ratio $\rho_F/\rho_E$ depends on $\Delta \phi_{qr}$,
so we cannot factor the exponential component out of the sum
and cancel the $\rho_{CN,\Delta \phi_{qr}}$ terms.
However, we can asymptotically expand $s_2/s_1$ for small
$B_0/c \coloneqq \epsilon \ll 1$;
additionally, since all of our simulations use $\beta \propto B$,
we also assume $\beta/c \sim \epsilon \ll 1$.
Finally, to simplify the calculation,
we define the asymptotic expansion of
$\rho_{NC,\Delta \phi_{qr}} \coloneqq
\rho_{NC,\Delta \phi_{qr}}^{(0)}
+
\rho_{NC,\Delta \phi_{qr}}^{(1)}
+
\order{\epsilon^2}
$
where
$
\rho_{NC,\Delta \phi_{qr}}^{(i)} \sim \epsilon^i
$
depends only on terms of total order $i$ in $B_0/c$ and $\beta/c$.

Then, to first order in $\epsilon$
we have
\begin{align*}
  \frac{s_2}{s_1}
  &=
  \frac
  {\sum_{r=1}^{m} \rho_{NC,\Delta \phi_{qr}}
    \exp \Bab{
      \delta (n-1)
      \bab{
        (n-1) c + n \beta(\Delta \phi_{qr}) (1 - 2 \alpha)
        - \frac{n-2}{2} B_0
      }
    }
  }
  {\sum_{r=1}^{m} \rho_{NC,\Delta \phi_{qr}}}
  \\
  &=
  e^{\delta (n-1)^2 c}
  \Biggl\{
  \\
  &\qquad
  \frac
  {\sum_{r=1}^{m}
    \pab{
      \rho_{NC,\Delta \phi_{qr}}^{(0)}
      +
      \rho_{NC,\Delta \phi_{qr}}^{(1)}
    }
    \Bab{1 +
      \delta (n-1)
      \bab{
        n \beta(\Delta \phi_{qr}) (1 - 2 \alpha)
        - \frac{n-2}{2} B_0
      }
    }
  }
  {\sum_{r=1}^{m} \pab{
    \rho_{NC,\Delta \phi_{qr}}^{(0)}
    +
    \rho_{NC,\Delta \phi_{qr}}^{(1)}
  }}
  \Biggr\}
  \\
  &\qquad
  + \order{\epsilon^2}
  \\
  &=
  e^{\delta (n-1)^2 c}
  \Biggl\{
  \frac
  {\sum_{r=1}^{m}
    \pab{
      \rho_{NC,\Delta \phi_{qr}}^{(0)}
      +
      \rho_{NC,\Delta \phi_{qr}}^{(1)}
    }
  }
  {\sum_{r=1}^{m} \pab{
    \rho_{NC,\Delta \phi_{qr}}^{(0)}
    +
    \rho_{NC,\Delta \phi_{qr}}^{(1)}
  }}
  \\
  &\qquad
  +
  \frac
  {\sum_{r=1}^{m} \rho_{NC,\Delta \phi_{qr}}^{(0)}
    \delta (n-1)
    \bab{
      n \beta(\Delta \phi_{qr}) (1 - 2 \alpha)
      - \frac{n-2}{2} B_0
    }
  }
  {\sum_{r=1}^{m} \rho_{NC,\Delta \phi_{qr}}^{(0)}
  }
  \Biggr\}
  + \order{\epsilon^2}
  \\
  &=
  e^{\delta (n-1)^2 c}
  \Bab{
    1
    - \delta (n-1) \frac{n-2}{2} B_0
    +
    \frac{\delta n (n-1)}{m} (1 - 2 \alpha)
    \sum_{r=1}^{m} \beta(\Delta \phi_{qr})
  }
  + \order{\epsilon^2}
  \\
  &=
  \exp \Bab{\delta (n-1) \bab{
    (n-1) c
    - \frac{n-2}{2} B_0
    +
    \frac{n (1 - 2 \alpha)}{m}
    \sum_{r=1}^{m} \beta(\Delta \phi_{qr})
  }}
  + \order{\epsilon^2}
  \\
  &=
  \exp \Bab{ \delta (n-1) \bab{
    (n-1) c
    - \frac{n-2}{2} B_0
    +
    \frac{n (1 - 2 \alpha)}{2} \beta_0
  }}
  + \order{\epsilon^2}
\end{align*}
Where, in the last line, we used
the definition of
$\beta(\Delta \phi_{qr}) = \beta_0 \bab{1 + \cos(\Delta \phi_{qr})}/2$
to write
$\sum_{r=1}^m \beta(\Delta \phi_{qr})
= \beta_0 \sum_{r=1}^m \Bab{1 + \cos\bab{2 \pi (q-r)/m}}/2
= m \beta_0/2$
since the cosine sum gives zero.

Finally, since each $\Delta \phi_i$ is equally likely,
the probability of fixing to \emph{any}
communicative state is $m s_1$;
likewise the total non-communicative fixation probability is $m s_2$.
Then, since the process must eventually absorb to either $C$ or $N$,
we have
$m s_1 + m s_2 = 1$, giving the probability of communicative fixation as
\begin{equation}
  m s_1 = \frac{1}{1 + s_2/s_1}
  \label{eq:full_analytic}
\end{equation}
with $s_2/s_1$ given by \cref{eq:full_analytic_frac},
and the asymptotic, small-$\epsilon$ approximation given by
\begin{equation}
  m s_1 = \frac{1}{1 + s_2/s_1}
  \approx
  \frac{1}
  {1 + \exp \Bab{\delta (n-1) \bab{
    (n-1) c
    - \frac{n-2}{2} B_0
    +
    \frac{n (1 - 2 \alpha)}{2} \beta_0
    }}
  }
  \label{eq:analytic_first_term}
\end{equation}


\section{Phase \vs{} Strategy}
\begin{figure*}
\begin{adjustwidth}{-2.25in}{0in}
  \centering
  \begin{subcaptiongroup}
    \stackinset{l}{20pt+\width-45pt+1.5cm}{t}{\width/240*207-45pt+1.5cm}%
      {\phantomcaption\label{fig:frac-comm-phase-fixed}\captiontext*{}}{%
    \stackinset{l}{20pt}{t}{\width/240*207-45pt+1.5cm}%
      {\phantomcaption\label{fig:frac-comm-baseline}\captiontext*{}}{%
    \stackinset{l}{20pt+\width-45pt+1.5cm}{t}{0pt}%
      {\phantomcaption\label{fig:chimera-strategy-fixed}\captiontext*{}}{%
    \stackinset{l}{20pt}{t}{0pt}%
      {\phantomcaption\label{fig:chimera-baseline}\captiontext*{}}{%
		{\includestandalone{tikz/c-elegans-fixed-aspect}}%
		}}}}
  \end{subcaptiongroup}
  \caption{
    \textbf{
			Comparison of strategy vs phase evolution.
    }
		Comparison of strategy vs phase evolution on the chimera-like index
		and communicative fraction.
    \protect{\subref{fig:chimera-baseline}},\protect{\subref{fig:chimera-strategy-fixed}}
		The chimera-like index is shown as a function of asymmetry $\alpha$.
    \protect{\subref{fig:chimera-baseline}}
		is the baseline case with both strategy and phase updated
		according to the birth-death Moran process,
		while
    \protect{\subref{fig:chimera-strategy-fixed}}
		keeps the $C/N$ strategy of each player fixed
		to its random initial value but subjects its phase $\phi_i$
		to the birth-death Moran process.
		The vertical error bars depict the standard deviation
		across ten seeds.
    \protect{\subref{fig:frac-comm-baseline}},\protect{\subref{fig:frac-comm-phase-fixed}}
		The communicative fraction $f_{\text{comm}}$ as a function
		of the maximum joint benefit $B_0$.
    \protect{\subref{fig:frac-comm-baseline}}
		is the baseline case with both strategy and phase updated
		according to the birth-death Moran process,
		while
    \protect{\subref{fig:frac-comm-phase-fixed}}
		has the phase $\phi_i$ of each player fixed
		to its random initial value while its $C/N$ strategy is subject
		to the birth-death Moran process.
  }\label{fig:phase-vs-strategy}
\end{adjustwidth}
\end{figure*}

We note that the communicative fraction is only sensitive to the
players' strategies, so fixing the strategies would be uninformative,
showing only the random initial distribution.
Likewise, the chimera-like index is only sensitive to the players'
phases,
so fixing the phases would be uninformative, showing only the random
initial distribution.

First, the \subref{fig:chimera-strategy-fixed} large chimera-like index
of the fixed-strategy case shows that phase-evolution is a key part of
chimera-like evolution.
Likewise, the \subref{fig:frac-comm-phase-fixed} similarity of the
fixed-phase case---and hence, large deviation from the well-mixed
case---demonstrates that strategy-evolution of \gls{celegans} alone
introduces qualitatively new
behavior compared to the well-mixed case.
Finally, the baseline, co-evolving strategy- and phase-dynamic case
differs from both the \subref{fig:chimera-strategy-fixed} fixed-strategy
and \subref{fig:frac-comm-phase-fixed} fixed-phase cases.
This implies that the co-evolving phase and strategy used in the main
text
introduces novel dynamics not captured by either the phase or strategy
alone.

\section{Data availability}
\subsection{Time-evolution animations}
The chimera-like quality of the games' populations
are necessarily time-dependent.
We found that animations showing the time-evolution
of players' strategies was a particularly useful way
to observe this chimera-like quality.
Therefore, we have included animations
showing the populations' evolving strategies
in the supplemental information for a subset of the parameters sampled.

\subsection{Interactive plots}
While discussing the simulation results in the main manuscript,
we necessarily focused on a particular region of parameter space.
Nevertheless, we also performed simulations across a range of parameters
to ensure the robustness of our results.
We used a Pluto.jl notebook to organize and visualize the data
and have made these results available for public analysis.
Additionally, this notebook includes time-evolution animations
for all of the simulated parameter combinations.

\bibliography{references.bib}

\end{document}
